\section{The single-degree-of-freedom oscillator}\label{sec:oscillator}

\subsection{Equation of motion}

A mass $m$ is connected to a moving base by a spring of stiffness $k_s$ and a viscous damper of
coefficient $c$. The base moves with displacement $y(t)$, so its acceleration is $\ddot y(t) =
x(t)$, the signal of interest. Let $u(t)$ be the absolute displacement of the mass and
\begin{equation}
  z(t) = y(t) - u(t)
\end{equation}
the \emph{relative displacement} (base minus mass). The spring and the damper act on the mass with
the force $k_s(y-u) + c(\dot y-\dot u) = k_s z + c\dot z$, so Newton's law reads $m\ddot u = k_s z +
c\dot z$. Since $u = y - z$, $\ddot u = x - \ddot z$, and
\begin{equation}
  m(x - \ddot z) = k_s z + c \dot z
  \quad\Longrightarrow\quad
  \ddot z + \frac{c}{m}\dot z + \frac{k_s}{m} z = x(t).
\end{equation}
With the \emph{natural angular frequency} $\wn = \sqrt{k_s/m} = 2\pi\fn$ and the \emph{damping
ratio} $\zeta = c/(2m\wn)$,
\begin{result}[Oscillator under base acceleration]
\begin{equation}\label{eq:sdof}
  \ddot z + 2\zeta\wn\dot z + \wn^2 z = x(t), \qquad Q = \frac{1}{2\zeta}.
\end{equation}
\end{result}
The quality factor $Q$ is the usual way to give the damping in vibration testing; $Q=10$ ($\zeta =
5\,\%$) is a common default. The stress in a component modelled by the oscillator is proportional to
the deformation of its spring, $S = K z$ with a constant $K$. The \emph{pseudo-acceleration}
$\wn^2 z$ is often plotted instead of $z$: at low frequency it equals the base acceleration.

\subsection{Frequency response}

Let the base acceleration be the complex exponential $x(t) = e^{\ii\omega t}$. The steady-state
response has the same time dependence, $z(t) = H(\omega)e^{\ii\omega t}$, with $\dot z = \ii\omega z$
and $\ddot z = -\omega^2 z$. Substituting in~\eqref{eq:sdof}:
\begin{equation}
  \bigl(-\omega^2 + 2\ii\zeta\wn\omega + \wn^2\bigr)H(\omega) = 1
  \quad\Longrightarrow\quad
  H(\omega) = \frac{1}{\wn^2-\omega^2 + 2\ii\zeta\wn\omega},
\end{equation}
and, since $|a+\ii b|^2 = a^2+b^2$,
\begin{result}[Frequency response]
\begin{equation}\label{eq:H2}
  |H(\omega)|^2 = \frac{1}{(\wn^2-\omega^2)^2 + (2\zeta\wn\omega)^2}.
\end{equation}
\end{result}
At $\omega=0$, $|H|=1/\wn^2$ (static deflection). At $\omega = \wn$ the first term of the
denominator vanishes and $|H(\wn)| = 1/(2\zeta\wn^2) = Q/\wn^2$: the resonance amplifies the static
response by $Q$.

\emph{Half-power bandwidth.} Near the resonance put $\omega = \wn(1+\varepsilon)$ with
$|\varepsilon|\ll1$. Then $\wn^2-\omega^2 = -\wn^2(2\varepsilon+\varepsilon^2) \approx -2\varepsilon\wn^2$
and $2\zeta\wn\omega\approx 2\zeta\wn^2$, so $|H|^2 \approx 1/\bigl(4\wn^4(\varepsilon^2+\zeta^2)\bigr)$.
This is half of the peak value $1/(4\zeta^2\wn^4)$ when $\varepsilon^2+\zeta^2 = 2\zeta^2$, i.e.\
$\varepsilon = \pm\zeta$. The band between the two half-power points has the width
$2\zeta\wn$, or in hertz
\begin{equation}\label{eq:bandwidth}
  \Delta f = 2\zeta\fn = \frac{\fn}{Q}.
\end{equation}
An oscillator with $Q=10$ responds mainly to a band of $\pm5\,\%$ around $\fn$, i.e.\ about $0.14$
octave wide. With 12 nodes per octave, adjacent oscillators ($2^{1/12}-1\approx6\,\%$ apart) share most
of their band. This overlap appears repeatedly in the design problems of Parts~II and~III.

\subsection{Response to a random input}\label{sec:respvar}

Let $x$ be a stationary random signal with PSD $G(f)$, represented by the lines of
Section~\ref{sec:psd}. Each line $A_m\cos(2\pi f_m t+\varphi_m)$ produces the response line
$|H(2\pi f_m)|A_m\cos(2\pi f_m t+\varphi_m+\arg H)$, so the response has the line powers
$|H|^2A_m^2/2$, and by~\eqref{eq:psddef} its PSD is
\begin{equation}\label{eq:Gz}
  G_z(f) = |H(2\pi f)|^2\, G(f).
\end{equation}
Its variance and, more generally, its \emph{spectral moments} are
\begin{equation}\label{eq:moments}
  m_k = \int_0^\infty f^k\, G_z(f)\,\dd f,\qquad k = 0,1,2,4,\qquad \sigma_z^2 = m_0 .
\end{equation}
Differentiation multiplies a line by $\ii\omega$, so the velocity $\dot z$ has the PSD
$(2\pi f)^2G_z$ and the variance $\sigma_{\dot z}^2 = (2\pi)^2 m_2$; likewise $\sigma_{\ddot z}^2 =
(2\pi)^4 m_4$. These moments describe how often the response crosses zero and how many peaks it has
(Section~\ref{sec:rice}).

\emph{White noise.} If $G(f)=G_0$ is constant, $\sigma_z^2 = G_0\int_0^\infty |H|^2\dd f$. Appendix~A
shows that $\int_0^\infty|H|^2\dd f = 1/(8\zeta\wn^3)$, hence
\begin{equation}\label{eq:miles}
  \sigma_z^2 = \frac{G_0}{8\zeta\wn^3} = \frac{G_0\,Q}{4\wn^3},
  \qquad
  \sigma_{\wn^2 z} = \sqrt{\wn^4\sigma_z^2} = \sqrt{\tfrac{\pi}{2}\,\fn\, Q\, G_0},
\end{equation}
the second form being Miles' equation for the RMS pseudo-acceleration (use $\wn = 2\pi\fn$:
$\wn^4 G_0Q/(4\wn^3) = 2\pi\fn G_0 Q/4 = \tfrac{\pi}{2}\fn QG_0$).

\subsection{The response variance is linear in the PSD: Lalanne's operator}\label{sec:lalanne}

In the program a PSD is described by its values $G_j$ at the nodes $f_{\mathrm{n},j}$ and is linear
between consecutive nodes. Because~\eqref{eq:moments} is linear in $G$, the response variance of the
oscillator $i$ is a linear combination of the nodal values,
\begin{result}[Lalanne's operator]
\begin{equation}\label{eq:W}
  \sigma_{z,i}^2 = \sum_{j=1}^n W_{ij}\,G_j ,
  \qquad W_{ij} = \int_0^\infty |H_i(f)|^2\,\phi_j(f)\,\dd f ,
\end{equation}
\end{result}
where $\phi_j$ is the ``hat'' function equal to 1 at $f_{\mathrm{n},j}$, 0 at the other nodes and
linear in between, and $H_i$ is the response of the oscillator $f_{\mathrm{n},i}$. Lalanne
integrated~\eqref{eq:W} in closed form; Appendix~\ref{app:lalanne} gives the complete derivation. The
matrix $W$ is computed once per damage model. Two details of the program matter:
\begin{itemize}[nosep]
\item the PSD of a synthesised profile is not cut at the first and last node but held constant over a
margin of $3/Q$ relative bandwidth beyond them (an oscillator at the edge of the grid would otherwise
see only half of its band); the methods of version~2.0 leave the effect of this margin to the closed
loop, while the operators of MMVSD include it (Section~\ref{sec:operators});
\item when the pre-filter of Section~\ref{sec:digital} is active, $|H_i|^2$ is multiplied by its power
response.
\end{itemize}

\subsection{How fast the response follows a change of the input}\label{sec:envelope}

Part~III needs to know how the response variance follows an input whose level changes with time.
The mechanical energy of the oscillator per unit mass is $E = \tfrac12\dot z^2 + \tfrac12\wn^2z^2$.
Its derivative, using~\eqref{eq:sdof} to replace $\ddot z$, is
\begin{equation}
  \frac{\dd E}{\dd t} = \dot z\ddot z + \wn^2 z\dot z = \dot z\,(\ddot z + \wn^2 z)
  = \dot z\,(x - 2\zeta\wn\dot z) = \dot z\,x - 2\zeta\wn\dot z^2 .
\end{equation}
Take expectations: $\dd\E[E]/\dd t = \E[\dot z x] - 2\zeta\wn\E[\dot z^2]$. Two facts close this
equation.

\emph{Equipartition.} In a lightly damped oscillator the response is a sinusoid at $\wn$ with a
slowly varying amplitude $A$; over a cycle $\dot z^2$ and $\wn^2z^2$ have the same mean $\wn^2A^2/2$,
so $\E[\dot z^2] = \wn^2\E[z^2] = \E[E]$.

\emph{Input power.} The term $\E[\dot z(t)x(t)]$ is the power that the base injects. The velocity is a
weighted sum of past inputs, $\dot z(t) = \int_0^\infty h_v(u)x(t-u)\,\dd u$ ($h_v$ is the velocity
response to a unit impulse of $x$), so $\E[\dot z(t)x(t)] = \int_0^\infty
h_v(u)\E[x(t-u)x(t)]\,\dd u$. For a broadband input the correlation $\E[x(t-u)x(t)]$ vanishes after a
time much shorter than one period of the oscillator: the injected power depends only on the current
intensity of the input near $\fn$, not on the energy already stored. Its value follows from the
stationary case, where $\dd\E[E]/\dd t = 0$: then $\E[\dot zx] = 2\zeta\wn\E[E] = 2\zeta\wn^3\sigma_z^2$,
and with~\eqref{eq:miles}, $\E[\dot z x] = 2\zeta\wn^3\,G_0/(8\zeta\wn^3) = G_0/4$.

Let now the input intensity near $\fn$ change slowly, $G_0\to G(t)$, and write $s(t) = \E[z^2] =
\E[E]/\wn^2$ for the \emph{local variance} of the response. Dividing the energy balance by $\wn^2$:
\begin{equation}
  \frac{\dd s}{\dd t} = \frac{G(t)}{4\wn^2} - 2\zeta\wn\, s
  = \frac{1}{\tau}\bigl(s_\infty(t) - s(t)\bigr),
  \qquad s_\infty(t) = \frac{G(t)}{8\zeta\wn^3},
\end{equation}
where $s_\infty(t)$ is the stationary variance for the current input~\eqref{eq:miles} and
\begin{result}[Envelope law of the response]
\begin{equation}\label{eq:tau}
  \tau = \frac{1}{2\zeta\wn} = \frac{Q}{2\pi\fn},\qquad
  s(t) = \int_0^\infty \frac{1}{\tau}e^{-u/\tau}\, s_\infty(t-u)\,\dd u .
\end{equation}
\end{result}
The second form is the solution of the first-order equation (differentiate it to check). The
response variance is the stationary variance low-pass filtered with the time constant $\tau$. For a
non-white input the same law is applied with $s_\infty(t) = \int|H|^2G(f,t)\dd f$ computed from the
local spectrum $G(f,t)$.

\subsection{Consequence: fast modulations do not reach the response}\label{sec:papoulis}

Suppose the stationary variance $s_\infty(t)$ fluctuates about its mean with variance $V$ and an
exponential correlation $\E[\delta s_\infty(t)\delta s_\infty(t+u)] = Ve^{-|u|/T_c}$, $T_c$ being the
correlation time (the duration of a burst). In the frequency domain the fluctuation has the two-sided
spectrum $2VT_c/(1+\omega^2T_c^2)$ (its integral over $\omega/2\pi$ is $V$), and the filter
of~\eqref{eq:tau} has $|H_\tau(\omega)|^2 = 1/(1+\omega^2\tau^2)$. The variance of the filtered
fluctuation is
\begin{equation}
  \Var[s] = \int_{-\infty}^{\infty}\frac{\dd\omega}{2\pi}\,
  \frac{2VT_c}{(1+\omega^2T_c^2)(1+\omega^2\tau^2)} = \frac{2VT_c}{2\pi}\cdot\frac{\pi}{T_c+\tau}
  = V\,\frac{T_c}{T_c+\tau},
\end{equation}
where the integral $\int_{-\infty}^{\infty}\dd\omega/\bigl((1+\omega^2a^2)(1+\omega^2b^2)\bigr) =
\pi/(a+b)$ follows from the partial fractions $\frac{1}{(1+\omega^2a^2)(1+\omega^2b^2)} =
\frac{1}{a^2-b^2}\bigl(\frac{a^2}{1+\omega^2a^2}-\frac{b^2}{1+\omega^2b^2}\bigr)$ and
$\int\dd\omega/(1+\omega^2a^2) = \pi/a$: $\frac{1}{a^2-b^2}(\pi a - \pi b) = \pi/(a+b)$.
The mean of $s$ is not changed by the filter (its gain at $\omega=0$ is 1). Inserting in the last
form of~\eqref{eq:kmod}, $\kappa = 3\tavg{s^2}/\tavg{s}^2 = 3(1+\Var[s]/\tavg{s}^2)$, the kurtosis of
the response is
\begin{result}[Low-pass effect of the oscillator on the kurtosis]
\begin{equation}\label{eq:papoulis}
  \kappa_z - 3 = \bigl(\kappa_\infty - 3\bigr)\,\frac{T_c}{T_c+\tau},
\end{equation}
\end{result}
where $\kappa_\infty-3 = 3V/\tavg{s}^2$ is the excess kurtosis that the response would have if it
followed the input instantaneously. Bursts much shorter than $\tau$ leave a nearly Gaussian response
(this is Papoulis' rule); bursts much longer than $\tau$ are fully transmitted. Since $\tau$ is
inversely proportional to $\fn$, the same input can produce Gaussian responses at low natural
frequencies and leptokurtic ones at high natural frequencies. The kurtosis of the response,
$\kappa_z(\fn)$, is therefore a function of the natural frequency, and it is this function, not the
kurtosis of $x$, that governs the damage (Section~\ref{sec:damage}).

\subsection{The oscillator as the program computes it}\label{sec:digital}

The program computes $z$ from the samples $x_k$ with a second-order recursion,
\begin{equation}\label{eq:recursion}
  z_k = b_1x_k + b_2x_{k-1} + b_3x_{k-2} - a_1 z_{k-1} - a_2 z_{k-2},
\end{equation}
whose coefficients are those of the \emph{ramp-invariant} discretisation: they give the exact
response of~\eqref{eq:sdof} at the sampling instants when $x$ varies linearly between samples.
Appendix~\ref{app:smallwood} derives them. The frequency response of the recursion is obtained by
inserting $x_k = e^{\ii\omega k\Delta t}$, $z_k = H_d e^{\ii\omega k\Delta t}$:
\begin{equation}\label{eq:Hd}
  H_d(f) = \frac{b_1 + b_2 e^{-\ii\theta} + b_3 e^{-2\ii\theta}}{1 + a_1e^{-\ii\theta}+a_2e^{-2\ii\theta}},
  \qquad \theta = 2\pi f\Delta t .
\end{equation}
It equals $H$ when $\fn\ll f_s$ but not when $\fn$ approaches $f_s/2$. The ratio of the response
variances to a flat input, $\int|H_d|^2/\int|H|^2$, is for $Q=10$:
\begin{center}
\begin{tabular}{lccccc}
\toprule
$f_s/\fn$ & 25 & 8.8 & 6.25 & 4.4 & 3.3 \\
variance ratio & 0.99 & 0.92 & 0.85 & 0.72 & 0.56 \\
\bottomrule
\end{tabular}
\end{center}
The FDS of a record and of a test are both computed with~\eqref{eq:recursion}, so the comparison is
consistent; but a design based on the analogue~\eqref{eq:H2} overestimates the response at high
natural frequencies when $f_s$ is only a few times $f_{\max}$. The methods of version~2.0 design with
$H$ and let the closed loop (Section~\ref{sec:closedloop}) correct the difference; MMVSD designs with
$H_d$ directly (Section~\ref{sec:operators}).

\emph{Pre-filter.} When $f_s\le 10f_{\max}$ the program first passes $x$ through a fixed third-order
digital filter, forward and then backward in time (zero phase), as in the FDS computation of the
dissertation. Applied twice, the filter multiplies the PSD by $|H_{\mathrm{pre}}(f)|^4$. The same
filter is applied to records and tests, and $|H_{\mathrm{pre}}|^4$ multiplies $|H|^2$ in every
operator of the program, so that the designs stay consistent with the measurement.
